\documentclass[10pt,a4paper]{amsart}
\usepackage{amsmath,amssymb,amsthm,graphicx}
% For commutative diagrams we use Xy-pic
\usepackage[all]{xy}

% We use 2cell for 2-commutative diagrams.
\xyoption{2cell}
% Theorem environments.
%
\theoremstyle{plain}
\newtheorem{theorem}[subsection]{Theorem}
\newtheorem{proposition}[subsection]{Proposition}
\newtheorem{lemma}[subsection]{Lemma}

\theoremstyle{definition}
\newtheorem{definition}[subsection]{Definition}
\newtheorem{example}[subsection]{Example}
\newtheorem{exercise}[subsection]{Exercise}
\newtheorem{situation}[subsection]{Situation}

\theoremstyle{remark}
\newtheorem{remark}[subsection]{Remark}
\newtheorem{remarks}[subsection]{Remarks}

\numberwithin{equation}{subsection}

% Macros
%
\def\lim{\mathop{\mathrm{lim}}\nolimits}
\def\colim{\mathop{\mathrm{colim}}\nolimits}
\def\Spec{\mathop{\mathrm{Spec}}}
\def\Hom{\mathop{\mathrm{Hom}}\nolimits}
\def\Ext{\mathop{\mathrm{Ext}}\nolimits}
\def\SheafHom{\mathop{\mathcal{H}\!\mathit{om}}\nolimits}
\def\SheafExt{\mathop{\mathcal{E}\!\mathit{xt}}\nolimits}
\def\Sch{\mathit{Sch}}
\def\Mor{\mathop{\mathrm{Mor}}\nolimits}
\def\Ob{\mathop{\mathrm{Ob}}\nolimits}
\def\Sh{\mathop{\mathit{Sh}}\nolimits}
\def\NL{\mathop{N\!L}\nolimits}
\def\CH{\mathop{\mathrm{CH}}\nolimits}
\def\proetale{{pro\text{-}\acute{e}tale}}
\def\etale{{\acute{e}tale}}
\def\QCoh{\mathit{QCoh}}
\def\Ker{\mathop{\mathrm{Ker}}}
\def\Im{\mathop{\mathrm{Im}}}
\def\Coker{\mathop{\mathrm{Coker}}}
\def\Coim{\mathop{\mathrm{Coim}}}

% Boxtimes
%
\DeclareMathSymbol{\boxtimes}{\mathbin}{AMSa}{"02}

%
% Macros for moduli stacks/spaces
%
\def\QCohstack{\mathcal{QC}\!\mathit{oh}}
\def\Cohstack{\mathcal{C}\!\mathit{oh}}
\def\Spacesstack{\mathcal{S}\!\mathit{paces}}
\def\Quotfunctor{\mathrm{Quot}}
\def\Hilbfunctor{\mathrm{Hilb}}
\def\Curvesstack{\mathcal{C}\!\mathit{urves}}
\def\Polarizedstack{\mathcal{P}\!\mathit{olarized}}
\def\Complexesstack{\mathcal{C}\!\mathit{omplexes}}
% \Pic is the operator that assigns to X its picard group, usage \Pic(X)
% \Picardstack_{X/B} denotes the Picard stack of X over B
% \Picardfunctor_{X/B} denotes the Picard functor of X over B
\def\Pic{\mathop{\mathrm{Pic}}\nolimits}
\def\Picardstack{\mathcal{P}\!\mathit{ic}}
\def\Picardfunctor{\mathrm{Pic}}
\def\Deformationcategory{\mathcal{D}\!\mathit{ef}}
\usepackage{cleveref}
\crefname{originallemma}{original lemma}{original lemmas}
\Crefname{originallemma}{Original lemma}{Original lemmas}
\makeatletter
\newlabel{coherent-lemma-chow-Noetherian}{{0200}{source}{}{sourcepointer-1}{}}
\newlabel{coherent-lemma-chow-Noetherian@cref}{{[originallemma][0][]0200}{[1][0][]source}{}{}{}}
\makeatother
\crefname{originallemma}{original lemma}{original lemmas}
\Crefname{originallemma}{Original lemma}{Original lemmas}
\makeatletter
\newlabel{coherent-remark-chow-Noetherian}{{0201}{source}{}{sourcepointer-2}{}}
\newlabel{coherent-remark-chow-Noetherian@cref}{{[originallemma][0][]0201}{[1][0][]source}{}{}{}}
\makeatother
\crefname{originallemma}{original lemma}{original lemmas}
\Crefname{originallemma}{Original lemma}{Original lemmas}
\makeatletter
\newlabel{divisors-lemma-quasi-projective-Noetherian-pic-effective-Cartier}{{0AYM}{source}{}{sourcepointer-3}{}}
\newlabel{divisors-lemma-quasi-projective-Noetherian-pic-effective-Cartier@cref}{{[originallemma][0][]0AYM}{[1][0][]source}{}{}{}}
\makeatother
\crefname{originallemma}{original lemma}{original lemmas}
\Crefname{originallemma}{Original lemma}{Original lemmas}
\makeatletter
\newlabel{lemma-c1-cap-additive}{{02SP}{source}{}{sourcepointer-4}{}}
\newlabel{lemma-c1-cap-additive@cref}{{[originallemma][0][]02SP}{[1][0][]source}{}{}{}}
\makeatother
\crefname{originallemma}{original lemma}{original lemmas}
\Crefname{originallemma}{Original lemma}{Original lemmas}
\makeatletter
\newlabel{lemma-geometric-cap}{{02SQ}{source}{}{sourcepointer-5}{}}
\newlabel{lemma-geometric-cap@cref}{{[originallemma][0][]02SQ}{[1][0][]source}{}{}{}}
\makeatother
\crefname{originallemma}{original lemma}{original lemmas}
\Crefname{originallemma}{Original lemma}{Original lemmas}
\makeatletter
\newlabel{lemma-pushforward-cap-c1}{{02SU}{source}{}{sourcepointer-6}{}}
\newlabel{lemma-pushforward-cap-c1@cref}{{[originallemma][0][]02SU}{[1][0][]source}{}{}{}}
\makeatother
\crefname{originallemma}{original lemma}{original lemmas}
\Crefname{originallemma}{Original lemma}{Original lemmas}
\makeatletter
\newlabel{varieties-definition-intersection-number}{{0BEP}{source}{}{sourcepointer-7}{}}
\newlabel{varieties-definition-intersection-number@cref}{{[originallemma][0][]0BEP}{[1][0][]source}{}{}{}}
\makeatother
\crefname{originallemma}{original lemma}{original lemmas}
\Crefname{originallemma}{Original lemma}{Original lemmas}
\makeatletter
\newlabel{varieties-lemma-chi-tensor-finite}{{0AYT}{source}{}{sourcepointer-8}{}}
\newlabel{varieties-lemma-chi-tensor-finite@cref}{{[originallemma][0][]0AYT}{[1][0][]source}{}{}{}}
\makeatother
\crefname{originallemma}{original lemma}{original lemmas}
\Crefname{originallemma}{Original lemma}{Original lemmas}
\makeatletter
\newlabel{varieties-lemma-intersection-number-additive}{{0BER}{source}{}{sourcepointer-9}{}}
\newlabel{varieties-lemma-intersection-number-additive@cref}{{[originallemma][0][]0BER}{[1][0][]source}{}{}{}}
\makeatother
\crefname{originallemma}{original lemma}{original lemmas}
\Crefname{originallemma}{Original lemma}{Original lemmas}
\makeatletter
\newlabel{varieties-lemma-intersection-number-and-pullback}{{0BET}{source}{}{sourcepointer-10}{}}
\newlabel{varieties-lemma-intersection-number-and-pullback@cref}{{[originallemma][0][]0BET}{[1][0][]source}{}{}{}}
\makeatother
\crefname{originallemma}{original lemma}{original lemmas}
\Crefname{originallemma}{Original lemma}{Original lemmas}
\makeatletter
\newlabel{varieties-lemma-numerical-intersection-effective-Cartier-divisor}{{0BEU}{source}{}{sourcepointer-11}{}}
\newlabel{varieties-lemma-numerical-intersection-effective-Cartier-divisor@cref}{{[originallemma][0][]0BEU}{[1][0][]source}{}{}{}}
\makeatother
\crefname{originallemma}{original lemma}{original lemmas}
\Crefname{originallemma}{Original lemma}{Original lemmas}
\makeatletter
\newlabel{varieties-lemma-numerical-polynomial-leading-term}{{0BEN}{source}{}{sourcepointer-12}{}}
\newlabel{varieties-lemma-numerical-polynomial-leading-term@cref}{{[originallemma][0][]0BEN}{[1][0][]source}{}{}{}}
\makeatother
\title{Numerical intersection numbers agree with the degree of the iterated first Chern classes capped with the fundamental class}
\author{The Stacks Project Authors}
\date{stacks-project:chow.tex:local-snapshot-20261012; GNU FDL 1.2 or later}
\begin{document}
\maketitle
\noindent\textit{Editorial note.} All delivered mathematical text is the author's own, copied from the pinned Stacks Project chapter file chow.tex; no statement or proof step was rewritten, reordered or removed.\par\medskip
\noindent\textit{Editorial note.} Preamble provenance: the pinned chapter file opens with \textbackslash{}input\{preamble\} and carries no preamble text of its own, so the project's single Stacks-wide preamble staged at sources/stacks-chow/preamble.tex is excerpted verbatim as two spans, lines 14-18 (Xy-pic setup) and lines 154-217 (the native theorem declarations and mathematical macros). Nothing was added, removed or redefined; the wrapper class and the amsmath/amssymb/amsthm/graphicx base packages are supplied by the assembler.\par\medskip
\noindent\textit{Editorial note.} Same-chapter pointers are located at their \textbackslash{}label lines inside the pinned chapter file; cross-chapter pointers are located in the project's own tag list staged in the snapshot (sources/stacks-chow/tags.txt), at the verbatim lines carrying their canonical tags.\par\medskip
\noindent\textit{Allowed standard prerequisites (editorial).}\begin{itemize}
\item Intersection numbers of invertible sheaves on a proper scheme. Varieties, labels varieties-definition-intersection-number (statement, proof line 8192) and varieties-lemma-chi-tensor-finite (proof line 8202) for the case d = 0. Supplies the left hand side being compared and its base case, and the decomposition of the intersection number over the components (proof line 8214).
\item Chow's lemma and the behaviour of intersection numbers under a birational proper pullback. Cohomology of Schemes, labels coherent-lemma-chow-Noetherian (proof line 8221) and coherent-remark-chow-Noetherian (proof line 8222); Varieties, label varieties-lemma-intersection-number-and-pullback (proof line 8227). Produces the H-projective birational model Y of X over which the computation may be done and shows the intersection number is unchanged by that pullback.
\item Compatibility of capping with a first Chern class with pushforward, additivity of both sides and the effective Cartier divisor case. Same chapter chow.tex, labels lemma-pushforward-cap-c1 (proof line 8233), lemma-c1-cap-additive (proof line 8244) and lemma-geometric-cap (proof line 8254); Divisors, label divisors-lemma-quasi-projective-Noetherian-pic-effective-Cartier (proof line 8241); Varieties, labels varieties-lemma-intersection-number-additive (proof line 8243) and varieties-lemma-numerical-intersection-effective-Cartier-divisor (proof line 8248). Moves the equality across the birational model, reduces an invertible sheaf to a difference of two effective Cartier divisors by additivity of both sides, and rewrites both sides in the divisor case so that the induction hypothesis applies.
\end{itemize}\medskip
\ifdefined\hypertarget\hypertarget{sourcepointer-1}{}\fi
\noindent\textit{Source reference.} Original lemma 0200: sources/stacks-chow/tags.txt, source lines 2617--2617. Navigation pointer for \textbackslash{}ref\{coherent-lemma-chow-Noetherian\} used inside the counted span (line 8221); the argument that consumes it is retained in the span.\par\medskip
\ifdefined\hypertarget\hypertarget{sourcepointer-2}{}\fi
\noindent\textit{Source reference.} Original lemma 0201: sources/stacks-chow/tags.txt, source lines 2618--2618. Navigation pointer for \textbackslash{}ref\{coherent-remark-chow-Noetherian\} used inside the counted span (line 8222); the argument that consumes it is retained in the span.\par\medskip
\ifdefined\hypertarget\hypertarget{sourcepointer-3}{}\fi
\noindent\textit{Source reference.} Original lemma 0AYM: sources/stacks-chow/tags.txt, source lines 13795--13795. Navigation pointer for \textbackslash{}ref\{divisors-lemma-quasi-projective-Noetherian-pic-effective-Cartier\} used inside the counted span (line 8241); the argument that consumes it is retained in the span.\par\medskip
\ifdefined\hypertarget\hypertarget{sourcepointer-4}{}\fi
\noindent\textit{Source reference.} Original lemma 02SP: sources/stacks-chow/pinned-original-source.tex, source lines 4330--4330. Navigation pointer for \textbackslash{}ref\{lemma-c1-cap-additive\} used inside the counted span (line 8244); the argument that consumes it is retained in the span.\par\medskip
\ifdefined\hypertarget\hypertarget{sourcepointer-5}{}\fi
\noindent\textit{Source reference.} Original lemma 02SQ: sources/stacks-chow/pinned-original-source.tex, source lines 4432--4432. Navigation pointer for \textbackslash{}ref\{lemma-geometric-cap\} used inside the counted span (line 8254); the argument that consumes it is retained in the span.\par\medskip
\ifdefined\hypertarget\hypertarget{sourcepointer-6}{}\fi
\noindent\textit{Source reference.} Original lemma 02SU: sources/stacks-chow/pinned-original-source.tex, source lines 4631--4631. Navigation pointer for \textbackslash{}ref\{lemma-pushforward-cap-c1\} used inside the counted span (line 8233); the argument that consumes it is retained in the span.\par\medskip
\ifdefined\hypertarget\hypertarget{sourcepointer-7}{}\fi
\noindent\textit{Source reference.} Original lemma 0BEP: sources/stacks-chow/tags.txt, source lines 14357--14357. Navigation pointer for \textbackslash{}ref\{varieties-definition-intersection-number\} used inside the counted span (line 8192); the argument that consumes it is retained in the span.\par\medskip
\ifdefined\hypertarget\hypertarget{sourcepointer-8}{}\fi
\noindent\textit{Source reference.} Original lemma 0AYT: sources/stacks-chow/tags.txt, source lines 13801--13801. Navigation pointer for \textbackslash{}ref\{varieties-lemma-chi-tensor-finite\} used inside the counted span (line 8202); the argument that consumes it is retained in the span.\par\medskip
\ifdefined\hypertarget\hypertarget{sourcepointer-9}{}\fi
\noindent\textit{Source reference.} Original lemma 0BER: sources/stacks-chow/tags.txt, source lines 14359--14359. Navigation pointer for \textbackslash{}ref\{varieties-lemma-intersection-number-additive\} used inside the counted span (line 8243); the argument that consumes it is retained in the span.\par\medskip
\ifdefined\hypertarget\hypertarget{sourcepointer-10}{}\fi
\noindent\textit{Source reference.} Original lemma 0BET: sources/stacks-chow/tags.txt, source lines 14361--14361. Navigation pointer for \textbackslash{}ref\{varieties-lemma-intersection-number-and-pullback\} used inside the counted span (line 8227); the argument that consumes it is retained in the span.\par\medskip
\ifdefined\hypertarget\hypertarget{sourcepointer-11}{}\fi
\noindent\textit{Source reference.} Original lemma 0BEU: sources/stacks-chow/tags.txt, source lines 14362--14362. Navigation pointer for \textbackslash{}ref\{varieties-lemma-numerical-intersection-effective-Cartier-divisor\} used inside the counted span (line 8248); the argument that consumes it is retained in the span.\par\medskip
\ifdefined\hypertarget\hypertarget{sourcepointer-12}{}\fi
\noindent\textit{Source reference.} Original lemma 0BEN: sources/stacks-chow/tags.txt, source lines 14356--14356. Navigation pointer for \textbackslash{}ref\{varieties-lemma-numerical-polynomial-leading-term\} used inside the counted span (line 8214); the argument that consumes it is retained in the span.\par\medskip
\begin{lemma}
\label{lemma-degrees-and-numerical-intersections}
Let $k$ be a field. Let $X$ be a proper scheme over $k$.
Let $Z \subset X$ be a closed subscheme of dimension $d$.
Let $\mathcal{L}_1, \ldots, \mathcal{L}_d$ be invertible
$\mathcal{O}_X$-modules. Then
$$
(\mathcal{L}_1 \cdots \mathcal{L}_d \cdot Z) =
\deg(
c_1(\mathcal{L}_1) \cap \ldots \cap c_1(\mathcal{L}_d) \cap [Z]_d)
$$
where the left hand side is defined in
Varieties, Definition \ref{varieties-definition-intersection-number}.
In particular,
$$
\deg_\mathcal{L}(Z) = \deg(c_1(\mathcal{L})^d \cap [Z]_d)
$$
if $\mathcal{L}$ is an ample invertible $\mathcal{O}_X$-module.
\end{lemma}

\begin{proof}
We will prove this by induction on $d$. If $d = 0$, then the result is
true by Varieties, Lemma \ref{varieties-lemma-chi-tensor-finite}.
Assume $d > 0$.

\medskip\noindent
Let $Z_i \subset Z$, $i = 1, \ldots, t$ be the irreducible components
of dimension $d$ with reduced induced scheme structure and let $m_i$ be the
multiplicity of $Z_i$ in $Z$. Then $[Z]_d = \sum m_i[Z_i]$ and
$c_1(\mathcal{L}_1) \cap \ldots \cap c_1(\mathcal{L}_d) \cap [Z]_d$
is the sum of the cycles
$m_i c_1(\mathcal{L}_1) \cap \ldots \cap c_1(\mathcal{L}_d) \cap [Z_i]$.
Since we have a similar decomposition for
$(\mathcal{L}_1 \cdots \mathcal{L}_d \cdot Z)$ by
Varieties, Lemma \ref{varieties-lemma-numerical-polynomial-leading-term}
it suffices to prove the lemma in case $Z = X$
is a proper variety of dimension $d$ over $k$.

\medskip\noindent
By Chow's lemma there exists a birational proper morphism $f : Y \to X$
with $Y$ H-projective over $k$. See Cohomology of Schemes, Lemma
\ref{coherent-lemma-chow-Noetherian} and Remark
\ref{coherent-remark-chow-Noetherian}. Then
$$
(f^*\mathcal{L}_1 \cdots f^*\mathcal{L}_d \cdot Y) =
(\mathcal{L}_1 \cdots \mathcal{L}_d \cdot X)
$$
by Varieties, Lemma \ref{varieties-lemma-intersection-number-and-pullback}
and we have
$$
f_*(c_1(f^*\mathcal{L}_1) \cap \ldots \cap c_1(f^*\mathcal{L}_d) \cap [Y]) =
c_1(\mathcal{L}_1) \cap \ldots \cap c_1(\mathcal{L}_d) \cap [X]
$$
by Lemma \ref{lemma-pushforward-cap-c1}. Thus we may replace $X$ by $Y$
and assume that $X$ is projective over $k$.

\medskip\noindent
If $X$ is a proper $d$-dimensional projective variety, then we can
write $\mathcal{L}_1 = \mathcal{O}_X(D) \otimes \mathcal{O}_X(D')^{\otimes -1}$
for some effective Cartier divisors $D, D' \subset X$
by Divisors, Lemma
\ref{divisors-lemma-quasi-projective-Noetherian-pic-effective-Cartier}.
By additivity for both sides of the equation
(Varieties, Lemma \ref{varieties-lemma-intersection-number-additive} and
Lemma \ref{lemma-c1-cap-additive})
we reduce to the case $\mathcal{L}_1 = \mathcal{O}_X(D)$ for some
effective Cartier divisor $D$.
By Varieties, Lemma
\ref{varieties-lemma-numerical-intersection-effective-Cartier-divisor}
we have
$$
(\mathcal{L}_1 \cdots \mathcal{L}_d \cdot X) =
(\mathcal{L}_2 \cdots \mathcal{L}_d \cdot D)
$$
and by Lemma \ref{lemma-geometric-cap} we have
$$
c_1(\mathcal{L}_1) \cap \ldots \cap c_1(\mathcal{L}_d) \cap [X] =
c_1(\mathcal{L}_2) \cap \ldots \cap c_1(\mathcal{L}_d) \cap [D]_{d - 1}
$$
Thus we obtain the result from our induction hypothesis.
\end{proof}

\end{document}
